\documentclass{article}
\usepackage{graphicx}

\title{A Lightweight Approach to Sentence Compression}
\author{Jane Doe \and Alex Kim}

\begin{document}
\maketitle

\begin{abstract}
We present a lightweight encoder-decoder for sentence compression that achieves
84.3 ROUGE-L on the Gigaword test set while using only 12M parameters. Our
approach matches the accuracy of a 350M-parameter baseline at one-thirtieth the
size, enabling on-device inference.
\end{abstract}

\section{Introduction}
Recent encoder-decoder models for summarization have grown to hundreds of
millions of parameters. We show that aggressive vocabulary tying and shared
embeddings preserve quality while reducing parameter count by a factor of 30
\cite{vaswani2017attention}.

\section{Method}
Our model uses six encoder and six decoder layers with shared input and output
embeddings. We follow the training recipe of \cite{vaswani2017attention} with
two adjustments described in Section~\ref{sec:training}.

\section{Training}
\label{sec:training}
We train for 120k steps on 8 V100 GPUs using a learning rate of 0.0005.

\section{Results}
On Gigaword our model achieves 84.3 ROUGE-L. See Table~\ref{tab:results}.

\begin{table}
\centering
\caption{ROUGE-L on Gigaword test set.}
\label{tab:results}
\begin{tabular}{lr}
Model & ROUGE-L \\
\hline
Baseline (350M) & 84.5 \\
Ours (12M) & 84.3 \\
\end{tabular}
\end{table}

\section{Conclusion}
A compact encoder-decoder can match a much larger baseline on sentence
compression at 84.3 ROUGE-L.

\bibliography{refs}
\bibliographystyle{plain}

\end{document}
